\documentclass[12pt,a4paper]{article}

\usepackage[T1]{fontenc}
\usepackage[utf8]{inputenc}
\usepackage[polish]{babel}
\usepackage{amsmath}
\usepackage[margin=2.5cm]{geometry}

\title{Nobliści fizyki i ich słynne równania}
\author{}
\date{}

\begin{document}

\maketitle

\begin{enumerate}
    \item \textbf{Max Planck} --- twórca teorii kwantów.
    \[
        E = h\nu
    \]
    Energia kwantu promieniowania jest proporcjonalna do jego częstotliwości.

    \item \textbf{Max Born} --- współtwórca mechaniki kwantowej.
    \[
        P(x,t) = \lvert\psi(x,t)\rvert^2
    \]
    Kwadrat modułu funkcji falowej opisuje prawdopodobieństwo znalezienia
    cząstki w danym miejscu i czasie.

    \item \textbf{Erwin Schrödinger} --- jeden z twórców mechaniki kwantowej.
    \[
        i\hbar\frac{\partial\psi}{\partial t}
        = \hat{H}\psi
    \]
    Jest to zależne od czasu równanie Schrödingera.

    \item \textbf{Albert Einstein} --- twórca szczególnej teorii względności.
    \[
        E = mc^2
    \]
    Energia i masa są równoważne, a współczynnikiem przeliczeniowym jest
    kwadrat prędkości światła.
\end{enumerate}

\end{document}